\documentclass{article}
\usepackage[T1]{fontenc}
\usepackage{lmodern}
\usepackage{graphicx}
\usepackage{amsmath,amssymb}
\usepackage[a4paper,margin=2cm]{geometry}
\usepackage[absolute,overlay]{textpos}
\setlength{\TPHorizModule}{1mm}
\setlength{\TPVertModule}{1mm}
\usepackage[indonesian]{babel}
\usepackage{hyperref}
\setlength{\emergencystretch}{2em}
% unit: Halaman 4

\begin{document}

% === Halaman 4 (python/native) ===

• Dalam teori algoritma, persoalan knapsack termasuk ke dalam 

kelompok NP-complete. 

• Persoalan yang termasuk NP-complete tidak dapat dipecahkan dalam 

orde waktu polinomial. 

• Ide dasar dari algoritma kriptografi knapsack adalah mengkodekan 

pesan sebagai rangkaian solusi dari persoalan knapsack. 

• Setiap bobot wi di dalam persoalan knapsack merupakan kunci rahasia, 

sedangkan bit-bit plainteks menyatakan bi. 

4

\end{document}
